\subsection{FORMAL POWER SERIES OVER A COMMUTATIVE RING}

For every set I, the additive monoid \(\mathbf{N}^{(I)}\) satisfies condition (D) of no. 10; for, if \(s = (n_i)_{i \in I}\) with \(n_i = 0\) except for the indices \(i\) in a finite subset \(\mathbf{H}\) of I, the relation \(s = t + u\) with \(t = (p_i)_{i \in I}\) and \(u = (q_i)_{i \in I}\) is equivalent to \(p_i + q_i = n_i\) for all \(i\) ; but this implies \(p_i = q_i = 0\) for \(i \notin \mathbf{H}\) and \(p_i \leqslant n_i\) , \(q_i \leqslant n_i\) for \(i \in \mathbf{H}\) ; there are therefore \(\prod_{i \in \mathbf{H}} (n_i + 1)\) ordered pairs \((t, u)\) in \(\mathbf{N}^{(I)}\) such that \(t + u = s\) .

We can therefore consider the total algebra of the monoid \(\mathbf{N}^{(I)}\) over A, which contains the (restricted) algebra \(\mathrm{A}[\mathbf{X}_i]_{i\in I}\) of this monoid. It is a unital commutative associative algebra called the algebra of formal power series in the indeterminates \(\mathbf{X}_i\) ( \(i\in I\) ) with coefficients in A and denoted by \(\mathrm{A}[[\mathbf{X}_i]]_{i\in I}\) ; its elements are called formal power series in the indeterminates \(\mathbf{X}_i\) ( \(i\in I\) ) with coefficients in A. Such an element \((\alpha_{\nu})_{\nu \in \mathbf{N}^{(I)}}\) is also denoted, following the convention made in no. 10, by \(\sum_{\nu \in \mathbf{N}^{(I)}}\alpha_{\nu}\mathbf{X}^{\nu}\) ; the \(\alpha_{\nu}\) are the coefficients of the formal power series and the \(\alpha_{\nu}\mathbf{X}^{\nu}\) its terms; a polynomial in the \(\mathbf{X}_i\) is therefore a formal power series with only a finite number of terms \(\neq 0\) .

Clearly an algebra isomorphism \(\mathbf{A}[[\mathbf{X}_i]]_{i\in I_1}\to \mathbf{A}[[\mathbf{X}_i]]_{i\in I_2}\) is canonically derived from every bijection \(\sigma \colon I_1\to I_2\) by mapping the formal power series \(\sum_{(n_i)}\alpha_{(n_i)}\cdot \prod_{i\in I_1}\mathbf{X}_i^{n_i}\) to the formal power series \(\sum_{(n_i)}\alpha_{(n_i)}\cdot \prod_{i\in I_1}\mathbf{X}_{\sigma (i)}^{n_i}\) .

Let \(J\) be a subset of \(I\) ; the algebra \(A[[X_i]]_{i \in J}\) can be identified with a subalgebra of \(A[[X_i]]_{i \in I}\) consisting of the formal power series \(\sum_{(n_i)} \alpha_{(n_i)} \cdot \prod_{i \in I} X_i^{n_i}\) where \(\alpha_{(n_i)} = 0\) for every element \((n_i) \in \mathbf{N}^{(I)}\) such that \(n_i \neq 0\) for at least one index \(i \in I - J\) . Further, if \(K = I - J\) , \(A[[X_i]]_{i \in I}\) is canonically identified with \((A[[X_j]]_{j \in J})[[X_k]]_{k \in K}\) , by identifying the formal power series \(\sum_{(n_i)} \alpha_{(n_i)} \cdot \prod_{i \in I} X_i^{n_i}\) with the formal power series \(\sum_{(m_k)} \beta_{(m_k)} \cdot \prod_{k \in K} X_k^{m_k}\) , where

\[
\beta_ {(m _ {k})} = \sum_ {(p _ {j})} \gamma_ {(p _ {j})} \cdot \prod_ {j \in J} X _ {j} ^ {p _ {j}}
\]

with \(\gamma_{(p_i)} = \alpha_{(n_i)}\) for the sequence \((n_i)\) such that \(n_i = p_i\) for \(i \in \mathbf{J}\) and \(n_i = m_i\) for \(i \in \mathbf{K}\) .

Given a formal power series \(u = \sum_{\nu} \alpha_{\nu} X^{\nu}\) , the terms \(\alpha_{\nu} X^{\nu}\) such that \(|\nu| = p\) are called the terms in \(u\) of total degree \(p\) . The formal power series \(u_p\) whose terms of total degree \(p\) are those of \(u\) and whose other terms are zero, is called the homogenous part of \(u\) of degree \(p\) ; when I is finite, \(u_p\) is a polynomial for all \(p\) ; \(u_0\) is identified with an element of A (also called the constant term of \(u\) ). If \(u\) and \(v\) are two formal power series and \(w = uv\) , then

\[
w _ {p} = \sum_ {r = 0} ^ {p} u _ {r} v _ {p - r}\tag{39}
\]

for every integer \(p \geqslant 0\) .

For every formal power series \(u \neq 0\) , the least integer \(p \geqslant 0\) such that \(u_{p} \neq 0\) is called the total order (or simply the order) of u. If this order is denoted by \(\omega(u)\) and u and v are two formal power series \(\neq 0\) , then

\[
\omega (u + v) \geqslant \inf (\omega (u), \omega (v)) \quad \text { if } u + v \neq 0\tag{40}
\]

\[
\omega (u v) \geqslant \omega (u) + \omega (v) \quad \text { if } u v \neq 0.\tag{41}
\]

Further, if \(\omega(u) \neq \omega(v)\) , then necessarily \(u + v \neq 0\) and the two sides of (40) are equal.

Note that the order of 0 is not defined. By an abuse of language, it is a convention to say that ``f is a formal power series of order \(\geqslant p\) (resp. \textgreater p)'' if the homogeneous part of f of degree n is zero for all n \textless{} p (resp. \(n \leqslant p\) ); 0 is therefore a ``formal power series of order \textgreater p'' for every integer \(p \geqslant 0\) .

Let \(J\) be a subset of \(I\) and let \(A[[X_i]]_{i \in I}\) be identified as above with \(B[[X_k]]_{k \in K}\) , where \(K = I - J\) and \(B = A[[X_j]]_{j \in J}\) ; corresponding to the above definitions applied to \(B[[X_k]]_{k \in K}\) there are new definitions for the formal power series \(u \in A[[X_i]]_{i \in I}\) ; a term \(\alpha_{(n_i)} \cdot \prod_{i \in I} X_i^{n_i}\) is said to be of degree \(p\) in the \(X_i\) of index \(i \in K\) if \(\sum_{i \in K} n_i = p\) and the formal power series of \(B[[X_k]]_{k \in K}\) with the same terms of degree \(p\) as \(u\) and the others zero is called the homogenous part of degree \(p\) in the \(X_i\) of index \(i \in K\) . If \(u \neq 0\) , the order \(\omega_K(u)\) with respect to the \(X_i\) of index \(i \in K\) is the smallest of the integers \(p \geqslant 0\) such that the homogeneous part of \(u\) of degree \(p\) in the \(X_i\) of index \(i \in K\) is \(\neq 0\) . Inequalities (40) and (41) still hold when \(\omega\) is replaced by \(\omega_K\) .

